\documentclass[11pt,a4paper]{article}
\usepackage[utf8]{inputenc}
\usepackage[T1]{fontenc}
\usepackage[brazil]{babel}
\usepackage[margin=2.5cm]{geometry}
\usepackage{titlesec}
\usepackage{enumitem}
\usepackage{booktabs}
\usepackage{tabularx}
\usepackage{array}
\usepackage{xcolor}
\usepackage{hyperref}
\usepackage{fancyhdr}

\definecolor{tfblue}{RGB}{20,40,70}
\definecolor{tfgold}{RGB}{200,150,40}

\titleformat{\section}{\normalfont\Large\bfseries\color{tfblue}}{\thesection}{0.6em}{}
\titleformat{\subsection}{\normalfont\large\bfseries\color{tfblue}}{\thesubsection}{0.6em}{}

\pagestyle{fancy}
\fancyhf{}
\fancyhead[L]{\small\textit{TrackFleet --- Escopo do Produto}}
\fancyhead[R]{\small\thepage}
\renewcommand{\headrulewidth}{0.4pt}

\newcolumntype{L}{>{\raggedright\arraybackslash}X}

\begin{document}

\begin{center}
{\LARGE\bfseries\color{tfblue} DOCUMENTO DE ESCOPO DO PROJETO}\\[2pt]
{\Large\bfseries TrackFleet --- Gestão de Rastreador de Automóveis}\\[4pt]
{\small\itshape Disciplina de Gerenciamento de Projetos $\cdot$ Framework PMBOK\textregistered\ 8\textsuperscript{a} Edição $\cdot$ Domínio: Escopo}\\[2pt]
{\small Integrantes: Enzo Allebrand e Kerlon Ribeiro (dois GPs) $\cdot$ 4\textsuperscript{o} semestre $\cdot$ Data: 20/08}
\end{center}

\vspace{2mm}
\noindent\textbf{\color{tfgold}\large Escopo do Produto}

\vspace{1mm}
\noindent\textit{O ``o quê'': a descrição das características, funções e atributos do resultado que o TrackFleet deve entregar. Este documento responde à pergunta ``o que a entrega deve ser e fazer?'' --- distinta do Escopo do Projeto (o ``como'') e da Qualidade, tratados em documentos separados.}

\section{Descrição do Produto}
O TrackFleet é uma plataforma web de rastreamento e gestão de frotas voltada a pequenas transportadoras. A plataforma integra-se, via API, a dispositivos de rastreamento GPS já existentes nos veículos da transportadora parceira (não há fabricação de hardware próprio) e oferece visibilidade quase em tempo real da frota, alertas de manutenção preventiva e indicadores de uso para apoiar decisões de gestão.

\section{Requisitos do Produto}

\subsection{Requisitos funcionais}
\begin{itemize}[itemsep=2pt]
  \item Cadastro de veículos, motoristas e vínculo entre eles.
  \item Integração com a API do provedor de rastreamento GPS piloto.
  \item Mapa com a localização quase em tempo real dos veículos ativos.
  \item Histórico de trajetos por veículo e por período.
  \item Geração automática de alertas de manutenção preventiva, com base em quilometragem rodada e/ou tempo decorrido.
  \item Painel de indicadores: veículos ativos, alertas pendentes, economia estimada.
  \item Autenticação e perfis de acesso distintos para gestor e motorista.
  \item Relatório de uso por veículo, exportável.
\end{itemize}

\subsection{Requisitos não funcionais}
\begin{itemize}[itemsep=2pt]
  \item Conformidade com a LGPD no tratamento de dados de localização e de motoristas (consentimento e minimização de dados).
  \item Interface responsiva, utilizável em desktop e navegador mobile.
  \item Disponibilidade da plataforma durante toda a janela de testes com a transportadora piloto.
\end{itemize}

\section{Fora do Escopo do Produto (Exclusões)}
Tão importante quanto listar o que entra é deixar explícito o que \textbf{não} entra, para evitar expectativas equivocadas:
\begin{itemize}[itemsep=2pt]
  \item Fabricação, aquisição ou instalação física de rastreadores GPS próprios.
  \item Aplicativo mobile nativo (iOS/Android) --- a primeira versão é web responsiva.
  \item Módulo de pagamento, faturamento ou cobrança.
  \item Roteirização ou otimização automática de rotas.
  \item Integração simultânea com múltiplos provedores de rastreamento (apenas um provedor piloto nesta fase).
\end{itemize}

\section{Estrutura do Escopo: Backlog do Produto}
Por se tratar de um projeto conduzido com governança híbrida e leve (conforme o Documento de Governança), o escopo do produto é estruturado como um \textbf{backlog}, decomposto em épicos, features e histórias de usuário, em vez de uma EAP rígida.

\begin{itemize}[itemsep=4pt]
  \item[\textbf{E1}] \textbf{Monitoramento em tempo real} --- localização atual e histórico de trajetos.
  \begin{itemize}
    \item Como gestor, quero ver todos os veículos ativos em um mapa, para saber onde a frota está agora.
    \item Como gestor, quero consultar o trajeto de um veículo em um período, para auditar o uso.
  \end{itemize}
  \item[\textbf{E2}] \textbf{Gestão de manutenção preventiva} --- alertas por km/tempo.
  \begin{itemize}
    \item Como gestor, quero receber um alerta antes que a manutenção vire corretiva, para reduzir custo.
  \end{itemize}
  \item[\textbf{E3}] \textbf{Gestão de frota e usuários} --- cadastros e perfis.
  \begin{itemize}
    \item Como gestor, quero cadastrar veículos e motoristas e vinculá-los, para organizar a frota.
  \end{itemize}
  \item[\textbf{E4}] \textbf{Relatórios e indicadores} --- painel e exportação.
  \begin{itemize}
    \item Como gestor, quero um painel com veículos ativos e alertas pendentes, para decidir rapidamente.
  \end{itemize}
\end{itemize}

\section{Value Breakdown Structure (VBS)}
A VBS conecta o escopo do produto ao valor pretendido, tornando visível qual entregável move mais o ponteiro do projeto.

\begin{center}
\begin{tabularx}{\textwidth}{L r}
\toprule
\textbf{Entregável (épico)} & \textbf{Valor} \\
\midrule
E1 --- Monitoramento em tempo real & 40\% \\
E2 --- Gestão de manutenção preventiva & 30\% \\
E4 --- Relatórios e indicadores & 20\% \\
E3 --- Gestão de frota e usuários (infraestrutura base) & 10\% \\
Conformidade com a LGPD & 100\% (obrigatório) \\
\bottomrule
\end{tabularx}
\end{center}

\noindent\textit{Leitura da VBS: o monitoramento em tempo real é o entregável mais valioso do produto e deve ser priorizado sobre relatórios avançados caso o prazo do semestre aperte. A conformidade com a LGPD é obrigatória e não é negociável, independentemente de prioridade de valor.}

\section{Critérios de Aceite do Produto}
O detalhamento dos critérios de aceite, da Definição de Pronto (DoD) e do processo de validação está no documento \textit{Qualidade}, que trata a qualidade como a terceira dimensão do escopo --- integrada, e não separada, das dimensões de produto e de projeto.

\end{document}
